

\subsection{Sequential Run}

To run the program using default settings from \texttt{.cfg} files:
\begin{lstlisting}[language=bash]
conda activate myenv  # if using anaconda
python main.py
\end{lstlisting}

\noindent 
The user may overwrite the default settings using arguments. e.g., the following command will make \texttt{P-GP-ID} to use \texttt{./data/VWT-data/simulation} as the simulation directory, and will overwrite the corresponding \texttt{.cfg} file.
\begin{lstlisting}[language=bash]
python main.py -p ./data/VWT-data/simulation
\end{lstlisting}

\noindent 
To see all available arguments:
\begin{lstlisting}[language=bash]
python main.py -h
\end{lstlisting}

\noindent Descriptions on each argument are in the \hyperref[SEC:args]{following section}.

\noindent Following is the list of default arguments:\\
\begin{tabular}{r | l | l}
  \hline \hline
   & Value & Meaning \\
  \hline
  \texttt{-p} & ./data/VWT-data/sphere & simulation directory \\
  \texttt{-m} &   & model name (\texttt{*.domain}) is automatically found \\
  \texttt{-A} & 0 180 1 & angular window will be set to [0, 180] DEG with 1 DEG interval \\
  \texttt{-S} & 0 & phi-sweep \\
  \texttt{-F} & 1000 2000 1 & frequency band will be set to [1000, 2000] MHz with 1 MHz interval \\
  \texttt{-v} &   & no validation \\
  \texttt{-a} & 0 & acquisition function: maximum variance \\
  \texttt{-x} & 1 & x-data normalization strategy: z-score (standard normalization) \\
  \texttt{-n} & 3 & number of initial frequency samples \\
  \texttt{--add} & 1 & added frequency samples per iteration \\
  \texttt{-t} & 3 & number of terms used for Gaussian process kernel \\
  \texttt{-s} & 1 & sampling strategy: Latin Hypercube Sampling \\
  \texttt{-i} & 20 & maximum iteration for Gaussian Process \\
  \texttt{-T} & 5e-5 & IVR tolerance; tolerance for Gaussian Process \\
  \hline \hline
\end{tabular}



\subsection{Parallel Run}

To run the program using MPI using 4 nodes:
\begin{lstlisting}[language=bash]
mpirun -n 4 python main.py
\end{lstlisting}

\noindent
When using \texttt{Anaconda}, the virtual environment needs to be activated on each computing node. The user may create a bash file:
\begin{lstlisting}[language=bash]
#!/bin/bash   file: run_node.sh
unset DISPLAY   # disables GNOME
export OMP_NUM_THREADS=$(nproc)   # allow each node to use maximum cores
source $(conda info --base)/etc/profile.d/conda.sh  # enables anaconda
conda activate myenv   # activates the virtual environment
python main.py   # runs P-GP-ID
\end{lstlisting}
and then run the bash code:
\begin{lstlisting}[language=bash]
mpirun -n 4 run_node.sh
\end{lstlisting}